
\documentclass{article}
\usepackage{colortbl}
\usepackage{makecell}
\usepackage{multirow}
\usepackage{supertabular}

\begin{document}

\newcounter{utterance}

\centering \large Interaction Transcript for game `codenames', experiment `abstract', episode 3 with reflexion{-}Qwen3.5{-}9B.
\vspace{24pt}

{ \footnotesize  \setcounter{utterance}{1}
\setlength{\tabcolsep}{0pt}
\begin{supertabular}{c@{$\;$}|p{.15\linewidth}@{}p{.15\linewidth}p{.15\linewidth}p{.15\linewidth}p{.15\linewidth}p{.15\linewidth}}
    \# & $\;$A & \multicolumn{4}{c}{Game Master} & $\;\:$B\\
    \hline

    \theutterance \stepcounter{utterance}  
    & & & \multicolumn{2}{p{0.3\linewidth}}{
        \cellcolor[rgb]{0.9,0.9,0.9}{
            \makecell[{{p{\linewidth}}}]{
                \texttt{\tiny{[GM$|$GM]}}
                {'hidden': {'team': ['endearment', 'randomization', 'suspicion'], 'innocent': ['amuck', 'avoidance', 'takings'], 'opponent': [], 'assassin': []}, 'revealed': {'team': {'team': [], 'innocent': [], 'opponent': [], 'assassin': []}, 'opponent': {'team': [], 'innocent': [], 'opponent': [], 'assassin': []}}}
            }
        }
    }
    & & \\ \\

    \theutterance \stepcounter{utterance}  
    & & \multicolumn{4}{p{0.6\linewidth}}{
        \cellcolor[rgb]{0.9,0.9,0.9}{
            \makecell[{{p{\linewidth}}}]{
                \texttt{\tiny{[P1$\langle$GM]}}
                \texttt{You have attempted this same task before and did not succeed. These are the plans you wrote after those attempts.} \\
\texttt{Trial \#1: New plan: I failed because my clues ("quality", "feeling", "state") were too broad and triggered the opponent's distractor words ("desirability", "despondency") instead of my targets. Next time, I will choose highly specific, unique semantic anchors that are unlikely to match any distractor words. For instance, instead of "state" for "impetuousness" and "damnation", I will use "reckless" (specific to impetuousness) or "sin" (specific to damnation), ensuring the clue is narrow enough to avoid accidental guesses of innocent words while still being recognizable to my partner.} \\
\\ 
\texttt{Trial \#2: New plan: I failed because I grouped semantically distinct words (like "endearment" and "suspicion") under a single clue ("doubt"), which caused my partner to guess an opponent's word ("viewpoint") or a distractor ("stigma") instead of my targets. Next time, I will strictly limit each clue to a maximum of two words that share a very specific, narrow semantic link (e.g., "endearment" and "value" under "affection") to minimize the risk of triggering opponent or distractor words. I will prioritize clearing high{-}risk words individually rather than batching them.} \\
\\ 
\texttt{Trial \#3: New plan: I failed because I used "doubt" as a clue for both "endearment" and "suspicion," creating a contradictory semantic link that likely confused my partner or triggered a distractor. "Doubt" is antonymous to "endearment," making it a poor anchor for that word. Next time, I will strictly separate words with opposing meanings into different clues. I will use a positive, affectionate anchor like "love" exclusively for "endearment" and a negative, investigative anchor like "suspicion" (or "fear") exclusively for "suspicion," ensuring each clue has a clear, unambiguous direction that avoids triggering distractors like "avoidance" or "amuck."} \\
\\ 
\\ 
\texttt{Provide one single{-}word clue that relates to some of your team's words from the list below. You can choose to how many and to which words your clue relates to.} \\
\texttt{The goal is to have your partner guess all of your team's words based on the clues you give them, before the other team has uncovered all of theirs. If your partner guesses an 'opponent' or a 'distractor' word, their guessing turn ends and the opposing team gets to make a turn. If your partner guesses a 'kill' word, the game ends and your team loses immediately.} \\
\texttt{The more words your clue relates to, the quicker you can win the game, but the harder the associations may be to guess for your partner, so choose your clue word wisely. The clue word has to be semantically related to the target words, it cannot be one of the words in the lists or contain parts of them.} \\
\\ 
\texttt{Always give your single{-}word clue and your comma{-}separated list of related target words in the following format and make your answers as short as possible, never include any other text than is required in this form:} \\
\\ 
\texttt{CLUE: <WORD>} \\
\texttt{TARGETS: <TARGETS>} \\
\\ 
\texttt{Your team words are: endearment, randomization, suspicion.} \\
\texttt{Your opponent's team words are: .} \\
\texttt{Distractor words are: amuck, avoidance, takings.} \\
\texttt{The kill words are: .} \\
            }
        }
    }
    & \\ \\

    \theutterance \stepcounter{utterance}  
    & \multicolumn{4}{p{0.6\linewidth}}{
        \cellcolor[rgb]{0.9,0.9,0.9}{
            \makecell[{{p{\linewidth}}}]{
                \texttt{\tiny{[P1$\rangle$GM]}}
                \texttt{CLUE: doubt} \\
\texttt{TARGETS: suspicion, randomization} \\
            }
        }
    }
    & & \\ \\

    \theutterance \stepcounter{utterance}  
    & & & \multicolumn{2}{p{0.3\linewidth}}{
        \cellcolor[rgb]{0.9,0.9,0.9}{
            \makecell[{{p{\linewidth}}}]{
                \texttt{\tiny{[GM$|$GM]}}
                \texttt{doubt} \\
            }
        }
    }
    & & \\ \\

    \theutterance \stepcounter{utterance}  
    & & & \multicolumn{2}{p{0.3\linewidth}}{
        \cellcolor[rgb]{0.9,0.9,0.9}{
            \makecell[{{p{\linewidth}}}]{
                \texttt{\tiny{[GM$|$GM]}}
                ['suspicion', 'randomization']
            }
        }
    }
    & & \\ \\

    \theutterance \stepcounter{utterance}  
    & & & \multicolumn{2}{p{0.3\linewidth}}{
        \cellcolor[rgb]{0.9,0.9,0.9}{
            \makecell[{{p{\linewidth}}}]{
                \texttt{\tiny{[GM$|$GM]}}
                {'word': 'suspicion', 'assignment': 'team'}
            }
        }
    }
    & & \\ \\

    \theutterance \stepcounter{utterance}  
    & & & \multicolumn{2}{p{0.3\linewidth}}{
        \cellcolor[rgb]{0.9,0.9,0.9}{
            \makecell[{{p{\linewidth}}}]{
                \texttt{\tiny{[GM$|$GM]}}
                {'word': 'randomization', 'assignment': 'team'}
            }
        }
    }
    & & \\ \\

    \theutterance \stepcounter{utterance}  
    & & & \multicolumn{2}{p{0.3\linewidth}}{
        \cellcolor[rgb]{0.9,0.9,0.9}{
            \makecell[{{p{\linewidth}}}]{
                \texttt{\tiny{[GM$|$GM]}}
                \texttt{opponent has won} \\
            }
        }
    }
    & & \\ \\

    \theutterance \stepcounter{utterance}  
    & \multicolumn{4}{p{0.6\linewidth}}{
        \cellcolor[rgb]{0.95,0.95,0.95}{
            \makecell[{{p{\linewidth}}}]{
                \texttt{\tiny{[P2$\rightarrow$P1]}}
                \texttt{Post{-}episode Reflexion (not sent to GM):} \\
\texttt{New plan: I failed because "doubt" is semantically contradictory to "endearment," creating a confusing link that likely triggered a distractor or opponent word. Next time, I will strictly separate words with opposing meanings into different clues. I will use a positive, affectionate anchor like "love" exclusively for "endearment" and a negative, investigative anchor like "fear" or "worry" exclusively for "suspicion," ensuring each clue has a clear, unambiguous direction that avoids triggering distractors like "avoidance" or "amuck." For "randomization," I will use a specific term like "shuffling" or "chance" in a separate clue to avoid broad associations.} \\
            }
        }
    }
    & & \\ \\

\end{supertabular}
}

\end{document}
